\section{Đặc tả giao thức}
\label{sec:giao-thuc}

Giao thức gồm hai giai đoạn: \textbf{handshake trên TCP} (trao đổi khóa và đăng nhập, xong thì đóng) và
\textbf{truyền dữ liệu trên UDP} (chat, DM, file, ACK, keep-alive). Mọi số nguyên trên dây đều là
\emph{big-endian}. Định nghĩa đầy đủ nằm trong \fname{internal/protocol} và \fname{docs/protocol.md}.

\subsection{Cổng}
\begin{table}[H]
  \centering
  \caption{Các cổng sử dụng}
  \label{tab:cong}
  \begin{tabularx}{\linewidth}{llY}
    \toprule
    \textbf{Cổng} & \textbf{Giao thức} & \textbf{Dùng cho} \\
    \midrule
    9000 & TCP & Handshake: trao đổi khóa, đăng nhập (đóng sau khi xong) \\
    9001 & UDP & Toàn bộ dữ liệu: message, file, ACK, PING, CLOSE \\
    9002 & UDP broadcast & Server quảng bá địa chỉ trong LAN để client tự tìm \\
    8080+ & TCP (HTTP tự viết) & Web UI của từng client, chỉ lắng nghe trên \texttt{127.0.0.1} \\
    \bottomrule
  \end{tabularx}
\end{table}

\subsection{Handshake TCP}
\label{subsec:handshake}
Mỗi message handshake có dạng \texttt{[4 byte độ dài][JSON]}. Độ dài bị giới hạn 1 MB: độ dài do bên kia tự
khai, nếu không giới hạn thì kẻ xấu chỉ cần gửi độ dài 4 GB là làm server cấp phát hết bộ nhớ.

\begin{table}[H]
  \centering
  \caption{Các message handshake}
  \label{tab:handshake}
  \begin{tabularx}{\linewidth}{clL{0.2\linewidth}Y}
    \toprule
    \textbf{\#} & \textbf{Hướng} & \textbf{Message} & \textbf{Nội dung} \\
    \midrule
    1 & C $\to$ S & \texttt{HELLO} & \texttt{client\_pub}: khóa công khai X25519 (tạm, mỗi phiên một khóa) \\
    2 & S $\to$ C & \texttt{KEYX} & \texttt{server\_pub}, \texttt{conn\_id} (định danh kết nối UDP),
    \texttt{udp\_port} \\
    3 & C $\to$ S & \texttt{AUTH} & \texttt{nonce} 12 byte ngẫu nhiên, \texttt{sealed} =
    AES-GCM$_{k_\text{tcp}}$(\{user, pass\}) \\
    4 & S $\to$ C & \texttt{AUTH\_OK} / \texttt{AUTH\_FAIL} & danh sách user online / lý do từ chối \\
    \bottomrule
  \end{tabularx}
\end{table}

Sau bước 2, hai bên tự tính cùng một bí mật chung mà không ai gửi nó qua mạng, rồi tách ra hai khóa:
\begin{align}
  \mathit{shared} &= \mathrm{X25519}(\mathit{priv}, \mathit{peerPub}) \\
  k_\text{tcp} &= \mathrm{SHA256}(\mathit{shared} \,\|\, \texttt{tcp}) \quad
    \text{mã hóa user/password trong \texttt{AUTH}} \\
  k_\text{udp} &= \mathrm{SHA256}(\mathit{shared} \,\|\, \texttt{udp}) \quad
    \text{mã hóa mọi gói UDP}
\end{align}
Dùng hai khóa riêng để hai kênh có cách sinh nonce khác nhau (TCP: nonce ngẫu nhiên; UDP: nonce từ header) không
bao giờ dùng chung một khóa. Sau \texttt{AUTH\_OK}, kết nối TCP đóng; client gửi gói UDP đầu tiên (PING) mang
\texttt{conn\_id} để server biết địa chỉ UDP của mình.

\subsection{Gói UDP}
\label{subsec:goi}

\begin{figure}[H]
  \centering
  \begin{tikzpicture}[x=0.82cm,y=1cm,font=\small]
    \foreach \xa/\xb/\fld/\sub in {0/2/{magic}/{0xC0DE},2/3/{ver}/{1},3/4/{flg}/{},4/6.5/{conn\_id}/{uint32},
      6.5/9/{packet\_num}/{uint32}} {
      \draw[fill=blue!8] (\xa,0) rectangle (\xb,1);
      \node[font=\footnotesize] at ({(\xa+\xb)/2},0.65) {\texttt{\fld}};
      \node[font=\scriptsize] at ({(\xa+\xb)/2},0.25) {\sub};
    }
    \draw[fill=orange!10] (9,0) rectangle (14.5,1);
    \node at (11.75,0.65) {frame (đã mã hóa)};
    \node[font=\scriptsize] at (11.75,0.25) {$\le$ 1172 byte};
    \draw[fill=gray!15] (14.5,0) rectangle (16,1);
    \node at (15.25,0.5) {tag};
    \foreach \x/\l in {0/0,2/2,3/3,4/4,6.5/8,9/12} \node[font=\scriptsize,above] at (\x,1) {\l};
    \draw[decorate,decoration={brace,amplitude=5pt,mirror}] (0,-0.1) -- (9,-0.1)
      node[midway,below=6pt,font=\scriptsize] {header 12 byte: gửi rõ, là AAD của AES-GCM};
    \draw[decorate,decoration={brace,amplitude=5pt,mirror}] (14.5,-0.1) -- (16,-0.1)
      node[midway,below=6pt,font=\scriptsize] {16 byte};
  \end{tikzpicture}
  \caption{Cấu trúc một gói UDP (tổng $\le$ 1200 byte)}
  \label{fig:goi}
\end{figure}

\begin{itemize}
  \item \textbf{Kích thước $\le$ 1200 byte} để không bị phân mảnh IP trên mọi đường truyền (MTU tối thiểu của
        IPv6 là 1280 byte, trừ header IP/UDP), giống QUIC \cite{rfc9000}. Phần frame còn lại
        $1200 - 12 - 16 = 1172$ byte.
  \item \texttt{magic} và \texttt{ver} giúp loại nhanh gói không thuộc giao thức trước khi tốn công giải mã.
  \item \texttt{conn\_id} cho biết gói thuộc kết nối nào: server dùng chung một socket UDP cho mọi client và
        tra \texttt{conn\_id} để chuyển gói tới đúng kết nối, nên client đổi địa chỉ (NAT) vẫn nhận ra.
  \item \texttt{packet\_num} \textbf{tăng mãi, không bao giờ dùng lại}, kể cả khi gửi lại: dữ liệu của gói bị mất
        được đóng vào gói mới với số mới. Nhờ vậy mỗi ACK chỉ khớp đúng một lần gửi, đo RTT không bao giờ mập mờ
        (không cần thuật toán Karn như TCP), và nonce mã hóa không bao giờ lặp (\cref{sec:bao-mat}).
\end{itemize}

\subsection{Frame}
\label{subsec:frame}
Phần đã mã hóa của gói chứa một hoặc nhiều frame xếp liền nhau. Bên nhận đọc byte \texttt{type} để biết độ dài
frame, cắt ra, rồi đọc tiếp tới hết. Nhờ vậy một gói có thể vừa mang dữ liệu vừa mang ACK (\emph{piggyback}).

\begin{table}[H]
  \centering
  \caption{Bốn loại frame}
  \label{tab:frame}
  \begin{tabularx}{\linewidth}{clL{0.3\linewidth}cY}
    \toprule
    \textbf{Type} & \textbf{Tên} & \textbf{Cấu trúc (byte)} & \textbf{ACK?} & \textbf{Ý nghĩa} \\
    \midrule
    0x01 & STREAM & \texttt{type(1) flags(1) stream\_id(4) offset(4) len(2) data} & Có &
    1 đoạn dữ liệu ($\le$ 1000 byte) của 1 stream; bit 0 của \texttt{flags} = FIN \\
    0x02 & ACK & \texttt{type(1) window(4) n(1)} $n \times$ \texttt{[start(4) end(4)]} & Không &
    SACK theo \texttt{packet\_num} ($n \le 8$) + cửa sổ nhận \\
    0x03 & PING & \texttt{type(1)} & Có & Giữ kết nối, báo địa chỉ client cho server \\
    0x04 & CLOSE & \texttt{type(1) reason\_len(1) reason} & Không & Đóng ngay, không chờ timeout \\
    \bottomrule
  \end{tabularx}
\end{table}

\paragraph{STREAM.} Một \emph{stream} là một dãy byte có thứ tự (một tin nhắn, hoặc một file). Nhiều stream
chạy song song trong cùng một kết nối. \texttt{offset} là vị trí byte đầu tiên của \texttt{data} trong stream,
bên nhận ghép lại theo offset. Client đánh \texttt{stream\_id} lẻ (1, 3, 5\dots), server đánh số chẵn, nên hai
bên mở stream đồng thời cũng không bao giờ trùng số. Mỗi frame mang tối đa 1000 byte dữ liệu để gói còn chỗ ghép
thêm một ACK (tối đa 70 byte).

\paragraph{ACK.} Liệt kê các khoảng \texttt{packet\_num} đã nhận, sắp giảm dần, khoảng đầu tiên kết thúc ở gói
lớn nhất. Ví dụ đã nhận gói 1--5, 7--9 và 12:
\begin{center}
  \texttt{[0x02][window=48000][n=3][12,12][7,9][1,5]} \quad (30 byte; thiếu 6, 10, 11)
\end{center}
Đây là \emph{selective ACK}: bên gửi biết chính xác gói nào thiếu và chỉ gửi lại dữ liệu của gói đó.
\texttt{window} là số byte bên nhận còn trống, dùng cho điều khiển luồng (\cref{subsec:dieu-toc}).
Gói chỉ chứa ACK thì không được ACK lại, nếu không hai bên sẽ ACK qua lại vô hạn.

\paragraph{CLOSE.} \texttt{reason} là chuỗi UTF-8 tối đa 255 byte; bên gửi tự cắt nếu dài hơn, và cắt đúng
ranh giới ký tự để không làm hỏng ký tự nhiều byte.

\subsection{Nội dung stream (tầng ứng dụng)}
\label{subsec:app}
Byte đầu tiên của mỗi stream cho biết loại:
\begin{center}
  \small
  \begin{tabular}{ll}
    MESSAGE: & \texttt{[0x01][JSON AppMsg \dots{} tới hết stream]} \\
    FILE: & \texttt{[0x02][meta\_len u32][JSON FileMeta][dữ liệu file \dots]}
  \end{tabular}
\end{center}

\begin{lstlisting}[caption={Cấu trúc nội dung tầng ứng dụng (\fname{appmsg.go}, rút gọn)},
  label=lst:appmsg]
type AppMsg struct {
	Type  string   `json:"type"` // chat | dm | users | join | leave | error
	From  string   `json:"from,omitempty"`
	To    string   `json:"to,omitempty"`
	Text  string   `json:"text,omitempty"`
	Users []string `json:"users,omitempty"`
	TS    int64    `json:"ts"`            // UnixMilli
	Seq   uint64   `json:"seq,omitempty"` // per-connection message number
}

type FileMeta struct {
	ID, From, To, Name string
	Size   int64
	SHA256 string // hex; receiver re-hashes the file to verify it
}
\end{lstlisting}

Meta của file đi \textbf{chung stream} với dữ liệu file thay vì gửi bằng một message riêng, vì giữa các stream
không có thứ tự: nếu tách ra thì dữ liệu file có thể tới trước meta. Bên nhận đọc meta theo từng byte tới
(\fname{FileHeaderParser}), không cần chờ cả file. SHA-256 của toàn bộ file đi trong meta để bên nhận kiểm tra file
nguyên vẹn đầu cuối, độc lập với tầng transport.

\paragraph{Thứ tự tin nhắn (\texttt{Seq}).} Mỗi tin là một stream riêng để tin này không phải chờ tin kia
(\cref{sec:hol}); hệ quả là khi mất gói, tin \#2 có thể giao trước tin \#1. Để vẫn giao đúng thứ tự gửi, mỗi bên
gửi đánh \texttt{Seq} = 1, 2, 3\dots{} theo từng kết nối, và bên nhận dùng \fname{protocol.SeqBuffer}: tin tới
sớm được giữ lại tới khi các tin trước đủ, tin trùng bị bỏ, \texttt{Seq} = 0 được giao ngay. Nếu giữ quá 4096 tin
(chỉ xảy ra khi có lỗi) thì bỏ qua chỗ trống để không treo vĩnh viễn. File không cần cơ chế này vì mỗi file là
một stream, tầng transport đã giao đúng thứ tự theo offset.

\phanB{Mô tả cách client và server dùng \texttt{Seq}/\texttt{SeqBuffer} (đánh số ở đâu, hàng đợi gửi của server
xử lý thế nào khi đầy) sau khi tích hợp xong.}

\subsection{State diagram}
\label{subsec:state}
\phanB{State diagram của client (và của session phía server): \textsc{Discovering} $\to$ \textsc{Handshake}
$\to$ \textsc{Connected} $\to$ \textsc{Reconnecting}/\textsc{Closed}. Trạng thái của kết nối transport
(\texttt{waiting}/\texttt{connected}/\texttt{closed}) có trong \fname{transport.Stats.State}.}
